% ============================================================================
% Guide theme: the shared visual identity for every guide's PDF.
%
% Loaded by pandoc after its own preamble (--include-in-header). Guide-specific
% values (title, category, colours, motif, trademarks, blurb) and build values
% (font path, edition, date) come from build/brand.tex, which build_guide_pdf.py
% generates from guide.toml and categories.json. See STYLE.md for the visual
% rationale and EDITORIAL.md for the house style the components serve.
% ============================================================================

% ── Page grid ─────────────────────────────────────────────────────────────────
% A 7.5 x 9.25 in trim, the size of many programming books, with generous
% margins so lines stay a comfortable 70-80 characters long.
\geometry{
  paperwidth=7.5in, paperheight=9.25in,
  top=0.95in, bottom=0.95in, left=0.85in, right=0.85in,
  headheight=14pt, headsep=18pt, footskip=30pt,
}

% ── Colour palette ────────────────────────────────────────────────────────────
% Neutrals are shared by every guide; Accent* come from brand.tex.
\definecolor{Ink}{HTML}{1E2329}        % body text
\definecolor{Muted}{HTML}{69707A}      % running heads, captions, labels
\definecolor{Faint}{HTML}{A7ADB5}      % hairlines on dark, placeholders
\definecolor{Rule}{HTML}{E3E0D8}       % hairlines, table rules, box borders
\definecolor{Panel}{HTML}{F7F6F2}      % code panels, table stripes
\definecolor{Night}{HTML}{16202E}      % cover, chapter bands, part pages
\definecolor{NightSoft}{HTML}{223044}  % motif strokes on Night
\definecolor{NoteInk}{HTML}{3B5B8C}    % Note callouts
\definecolor{NoteTint}{HTML}{EEF2F8}
\definecolor{WarnInk}{HTML}{B45309}    % Warning callouts
\definecolor{WarnTint}{HTML}{FDF4E6}
\color{Ink}

% ── Typography ────────────────────────────────────────────────────────────────
% Source Serif 4 for reading, Inter for structure (headings, labels, tables),
% JetBrains Mono for code. All three are SIL Open Font License fonts shipped
% with the guide, so every build looks identical.
\setmainfont{SourceSerif4}[
  Path=\GuideFontPath, Extension=.otf,
  UprightFont=*-Regular, ItalicFont=*-It,
  BoldFont=*-Semibold, BoldItalicFont=*-SemiboldIt,
  Scale=1.0, Ligatures=TeX, Numbers=OldStyle]
\setsansfont{Inter}[
  Path=\GuideFontPath, Extension=.otf,
  UprightFont=*-Regular, ItalicFont=*-Italic,
  BoldFont=*-SemiBold, BoldItalicFont=*-Italic,
  Scale=MatchLowercase, Ligatures=TeX]
\setmonofont{JetBrainsMono}[
  Path=\GuideFontPath, Extension=.ttf,
  UprightFont=*-Regular, ItalicFont=*-Italic,
  BoldFont=*-Bold, BoldItalicFont=*-BoldItalic,
  Scale=MatchLowercase, HyphenChar=None,
  RawFeature={-calt;-liga}]% no ligatures (=> stays =>) and no hyphens in code
\newfontfamily\HeadFont{Inter}[Path=\GuideFontPath, Extension=.otf,
  UprightFont=*-Bold, BoldFont=*-ExtraBold, Scale=MatchLowercase, Ligatures=TeX]
\newfontfamily\HeadSemi{Inter}[Path=\GuideFontPath, Extension=.otf,
  UprightFont=*-SemiBold, BoldFont=*-Bold, Scale=MatchLowercase, Ligatures=TeX]
\newfontfamily\LabelFont{Inter}[Path=\GuideFontPath, Extension=.otf,
  UprightFont=*-Medium, BoldFont=*-SemiBold, Scale=MatchLowercase,
  LetterSpace=6, Ligatures=TeX]
\newfontfamily\DisplayFont{Inter}[Path=\GuideFontPath, Extension=.otf,
  UprightFont=*-ExtraBold, BoldFont=*-Black, Scale=MatchLowercase, WordSpace={1.9,1,1},
  Ligatures=TeX]
% A symbol fallback for characters the code font lacks, such as the test
% runner's check marks (used by guide.lua through \GuideSym).
\newfontfamily\GuideSymbolFont{DejaVuSans}[Path=\GuideFontPath, Extension=.ttf,
  Scale=MatchLowercase]
\newcommand{\GuideSym}[1]{{\GuideSymbolFont #1}}
\newfontfamily\LiningSerif{SourceSerif4}[Path=\GuideFontPath, Extension=.otf,
  UprightFont=*-Regular, ItalicFont=*-It, BoldFont=*-Semibold, Numbers=Lining,
  Ligatures=TeX]

% ── Script fonts (opt-in) ─────────────────────────────────────────────────────
% A guide that teaches a language with its own script (guide.toml:
% language = "ja") gets that script's fonts through xeCJK: Noto Serif JP beside
% Source Serif 4 in the body, Noto Sans JP beside Inter in headings, labels and
% tables. brand.tex defines \GuideCJKFontPath only for such guides, so every
% other guide builds exactly as before.
\ifdefined\GuideCJKFontPath
  \usepackage{xeCJK}
  \setCJKmainfont{\GuideCJKSerif}[Path=\GuideCJKFontPath, Extension=.otf,
    UprightFont=*-Regular, BoldFont=*-Bold, ItalicFont=*-Regular,
    BoldItalicFont=*-Bold]
  \setCJKsansfont{\GuideCJKSans}[Path=\GuideCJKFontPath, Extension=.otf,
    UprightFont=*-Regular, BoldFont=*-Bold, ItalicFont=*-Regular,
    BoldItalicFont=*-Bold]
  \setCJKmonofont{\GuideCJKSans}[Path=\GuideCJKFontPath, Extension=.otf,
    UprightFont=*-Regular, BoldFont=*-Bold]
  \setCJKfamilyfont{GuideCJKHead}{\GuideCJKSans}[Path=\GuideCJKFontPath,
    Extension=.otf, UprightFont=*-Bold, BoldFont=*-Bold]
  \setCJKfamilyfont{GuideCJKLabel}{\GuideCJKSans}[Path=\GuideCJKFontPath,
    Extension=.otf, UprightFont=*-Regular, BoldFont=*-Bold]
  % Headings and labels set their Latin face with \HeadFont and friends; give
  % each the matching Japanese sans so a heading never mixes in the serif.
  \let\GuideLatinHeadFont\HeadFont
  \renewcommand{\HeadFont}{\GuideLatinHeadFont\CJKfamily{GuideCJKHead}}
  \let\GuideLatinHeadSemi\HeadSemi
  \renewcommand{\HeadSemi}{\GuideLatinHeadSemi\CJKfamily{GuideCJKHead}}
  \let\GuideLatinLabelFont\LabelFont
  \renewcommand{\LabelFont}{\GuideLatinLabelFont\CJKfamily{GuideCJKLabel}}
  \let\GuideLatinDisplayFont\DisplayFont
  \renewcommand{\DisplayFont}{\GuideLatinDisplayFont\CJKfamily{GuideCJKHead}}
  % Western punctuation stays with the Western text around it; xeCJK would
  % otherwise set curly quotes, dashes and ellipses full width ("Let’ s").
  \xeCJKDeclareCharClass{Default}{"00B7, "2013, "2014, "2018, "2019, "201C,
    "201D, "2026}
  % Box drawing and math symbols appear in kaomoji (┻━┻, (ﾟ∀ﾟ)); the Latin
  % fonts lack them and the Japanese fonts have them.
  \xeCJKDeclareCharClass{CJK}{"2200 -> "22FF, "2500 -> "257F}
\fi

\linespread{1.18}
\setlength{\parindent}{0pt}
\setlength{\parskip}{0.62em plus 2pt minus 1pt}
\setlength{\emergencystretch}{3em}
\raggedbottom
\widowpenalty=10000
\clubpenalty=10000

% Inline code: monospaced in a deeper shade of the accent, so it reads as code
% without a box that would break across lines.
\let\GuideOrigTexttt\texttt
% (\leavevmode stops a colour change at the start of a table cell or list item
% from opening an empty line.)
\renewcommand{\texttt}[1]{\leavevmode{\color{AccentDark}\GuideOrigTexttt{#1}}}

% Links: accent colour, no boxes. Internal and external links look the same;
% readers of a book should not need to know the difference.
% (pandoc loads hyperref after this file, so apply the settings at \begin{document}.)
\AtBeginDocument{%
  \hypersetup{colorlinks=true, linkcolor=AccentDark, urlcolor=AccentDark,
    citecolor=AccentDark, filecolor=AccentDark,
    pdfauthor={}, pdfcreator={}, pdfsubject={\GuideSubtitle}}%
  \urlstyle{same}}
\PassOptionsToPackage{hypertexnames=false}{hyperref}

% Lists: accent bullets, a little air, aligned with the text block.
\setlist{itemsep=0.25em, topsep=0.35em, parsep=0.15em, leftmargin=1.5em}
\setlist[itemize,1]{label={\color{Accent}\small\textbullet}}
\setlist[itemize,2]{label={\color{Muted}\scriptsize\textendash}}
\setlist[enumerate,1]{label={\sffamily\small\bfseries\color{AccentDark}\arabic*.}}
\renewcommand{\tightlist}{\setlength{\itemsep}{0.18em}\setlength{\parskip}{0pt}}

% Figures and diagrams: never wider than the text or taller than about half a
% page, so a figure and its caption always fit on one page. (pandoc 3.x sizes
% images with \pandocbounded, which would otherwise allow a full page.)
\makeatletter
\def\maxheight{\ifdim\Gin@nat@height>0.5\textheight 0.5\textheight\else\Gin@nat@height\fi}
\newsavebox\GuideImageBox
\AtBeginDocument{%
  % pandoc defines \pandocbounded only for books with images; define it here so
  % a book without any (the Japanese guide) still compiles.
  \providecommand*\pandocbounded[1]{#1}%
  \renewcommand*\pandocbounded[1]{%
    \sbox\GuideImageBox{#1}%
    \Gscale@div\@tempa{0.55\textheight}{\dimexpr\ht\GuideImageBox+\dp\GuideImageBox\relax}%
    \Gscale@div\@tempb{\linewidth}{\wd\GuideImageBox}%
    \ifdim\@tempb\p@<\@tempa\p@\let\@tempa\@tempb\fi
    \ifdim\@tempa\p@<\p@\scalebox{\@tempa}{\usebox\GuideImageBox}%
    \else\usebox{\GuideImageBox}\fi}}
\makeatother

% ── Section headings (inside a lesson) ────────────────────────────────────────
\titleformat{\section}
  {\HeadFont\fontsize{15}{19}\selectfont\color{Ink}}{}{0pt}{}
\titlespacing*{\section}{0pt}{1.9em plus 0.4em}{0.5em}
\titleformat{\subsection}
  {\HeadSemi\fontsize{11.5}{15}\selectfont\color{AccentDark}}{}{0pt}{}
\titlespacing*{\subsection}{0pt}{1.4em plus 0.3em}{0.3em}
\titleformat{\subsubsection}
  {\HeadSemi\normalsize\color{Ink}}{}{0pt}{}
\titlespacing*{\subsubsection}{0pt}{1.1em}{0.2em}
\titleformat{\paragraph}[runin]
  {\HeadSemi\normalsize\color{Ink}}{}{0pt}{}[.]

% "Going deeper" subsections get a small label above the title.
\newcommand{\GuideDeeper}[1]{%
  \par\vspace{1.2em}\needspace{5\baselineskip}%
  {\LabelFont\scriptsize\color{Muted}GOING DEEPER}\par\vspace{0.1em}%
  {\HeadSemi\fontsize{11.5}{15}\selectfont\color{AccentDark}#1}\par\vspace{0.25em}}

% ── Running heads and folios ──────────────────────────────────────────────────
\newcommand{\GuideChapterMark}{}
\newcommand{\GuideLessonMark}{}
\fancypagestyle{guide}{%
  \fancyhf{}%
  \fancyhead[L]{\LabelFont\scriptsize\color{Muted}\MakeUppercase{\GuideChapterMark}}%
  \fancyhead[R]{\sffamily\scriptsize\color{Muted}\GuideLessonMark}%
  \fancyfoot[L]{\sffamily\scriptsize\color{Faint}\GuideTitle}%
  \fancyfoot[R]{\HeadSemi\small\color{AccentDark}\thepage}%
  \renewcommand{\headrulewidth}{0.4pt}%
  \renewcommand{\headrule}{{\color{Rule}\hrule height\headrulewidth width\headwidth}}%
  \renewcommand{\footrulewidth}{0pt}}
\fancypagestyle{opener}{%
  \fancyhf{}%
  \fancyfoot[R]{\HeadSemi\small\color{AccentDark}\thepage}%
  \renewcommand{\headrulewidth}{0pt}\renewcommand{\footrulewidth}{0pt}}
\fancypagestyle{plain}{% LaTeX's own \chapter and \part pages use this
  \fancyhf{}\fancyfoot[R]{\HeadSemi\small\color{AccentDark}\thepage}%
  \renewcommand{\headrulewidth}{0pt}}
\pagestyle{guide}

% ── Motifs (cover, chapter bands, part pages) ─────────────────────────────────
% Each guide picks one in brand.tex; the geometry is shared.
\tikzset{
  motif hex/.pic={\draw[line width=0.7pt] (0:1) \foreach \a in {60,120,...,359} {-- (\a:1)} -- cycle;},
}
\newcommand{\GuideMotifHexagons}[4]{% colour, x0, y0, rows
  \foreach \r in {0,...,#4} {\foreach \c in {0,...,9} {
    \pgfmathsetmacro{\x}{#2 + \c*1.5}\pgfmathsetmacro{\y}{#3 + \r*1.732 + mod(\c,2)*0.866}
    \path[draw=#1] (\x,\y) pic[scale=0.82] {motif hex};}}}
\newcommand{\GuideMotifDots}[4]{%
  \foreach \r in {0,...,#4} {\foreach \c in {0,...,22} {
    \fill[#1] (#2 + \c*0.6, #3 + \r*0.6) circle (0.05);}}}
% A few highlighted cells on the same grid (column, row pairs), for the accent.
\newcommand{\GuideMotifHexCells}[2]{% colour, list of c/r
  \foreach \c/\r in {#2} {
    \pgfmathsetmacro{\x}{\c*1.5}\pgfmathsetmacro{\y}{\r*1.732 + mod(\c,2)*0.866}
    \path[draw=#1, line width=1pt] (\x,\y) pic[scale=0.82] {motif hex};}}
\newcommand{\GuideMotifDotCells}[2]{%
  \foreach \c/\r in {#2} {\fill[#1] (\c*0.6, \r*0.6) circle (0.09);}}
\newcommand{\GuideMotifAccent}[2]{%
  \ifdefstring{\GuideMotifName}{hexagons}{\GuideMotifHexCells{#1}{#2}}{\GuideMotifDotCells{#1}{#2}}}
\newcommand{\GuideMotif}[4]{%
  \ifdefstring{\GuideMotifName}{hexagons}{\GuideMotifHexagons{#1}{#2}{#3}{#4}}{\GuideMotifDots{#1}{#2}{#3}{#4}}}

% ── Cover and colophon ────────────────────────────────────────────────────────
\newcommand{\GuideCover}{%
  \begin{titlepage}
  \pagenumbering{alph}% keep cover page labels apart from the numbered pages
  \begin{tikzpicture}[remember picture, overlay]
    \fill[Night] (current page.south west) rectangle (current page.north east);
    \begin{scope}[shift={(current page.south east)}, xshift=-11cm, yshift=-0.5cm]
      \clip (0,0) rectangle (11.5cm, 11cm);
      \GuideMotif{NightSoft}{0}{0}{6}
      \GuideMotifAccent{AccentBright}{4/1, 5/1, 5/2}
    \end{scope}
    \fill[AccentBright] ([xshift=0.85in, yshift=-1.55in]current page.north west) rectangle ++(1.1in, 5pt);
    \node[anchor=north west, text width=5.6in, inner sep=0pt, align=left]
      at ([xshift=0.85in, yshift=-0.95in]current page.north west)
      {\LabelFont\small\color{AccentBright}\MakeUppercase{\GuideCategory}};
    \node[anchor=north west, text width=5.8in, inner sep=0pt, align=left]
      at ([xshift=0.85in, yshift=-1.9in]current page.north west)
      {\hyphenpenalty=10000\exhyphenpenalty=10000\raggedright
       {\HeadFont\bfseries\fontsize{46}{52}\selectfont\spaceskip=0.26em\relax\color{white}\GuideTitle\par}
       \vspace{18pt}
       {\LiningSerif\itshape\fontsize{15}{21}\selectfont\color{Faint}\GuideSubtitle\par}};
    \node[anchor=south west, inner sep=0pt, align=left, text width=4.5in]
      at ([xshift=0.85in, yshift=0.8in]current page.south west)
      {\setlength{\parskip}{0pt}%
       {\HeadSemi\small\color{white}\GuideCoverage\par}\vspace{5pt}%
       {\LabelFont\scriptsize\color{Faint}\MakeUppercase{\GuideEditionLine}\par}};
  \end{tikzpicture}
  \end{titlepage}
  % Copyright page on the back of the cover.
  \thispagestyle{empty}
  \null\vfill
  {\sffamily\scriptsize\color{Muted}\setlength{\parskip}{0.8em}\raggedright
   {\HeadSemi\small\color{Ink}\GuideTitle}\par\vspace{-0.6em}
   \GuideSubtitle\par
   \GuideEditionLine. \GuideCoverage.\par
   \GuideCopyright\par
   \GuideTrademarks\par
   The information in this book is distributed on an ``as is'' basis, without
   warranty. Every effort has been made to ensure its accuracy, but neither the
   publisher nor the authors accept liability for any loss or damage caused, or
   alleged to have been caused, directly or indirectly by the instructions or
   code it contains.\par
   Set in Source Serif 4, Inter and JetBrains Mono. Built \GuideBuildDate.\par}
  \clearpage
  \pagenumbering{roman}
  \pagestyle{opener}}
\renewcommand{\maketitle}{\GuideCover}

% ── Contents ──────────────────────────────────────────────────────────────────
\renewcommand{\contentsname}{Contents}
\titlecontents{part}[2.3em]{\addvspace{0.95em}\HeadSemi\normalsize\color{Ink}}
  {\contentslabel[{\color{Accent}\thecontentslabel}]{2.3em}}{\hspace*{-2.3em}}
  {\hfill\color{AccentDark}\contentspage}[\vspace{0.1em}]
\titlecontents{chapter}[4.9em]{\addvspace{0.05em}\LiningSerif\small\color{Ink}}
  {\contentslabel[{\sffamily\footnotesize\color{Muted}\thecontentslabel}]{2.6em}}{\hspace*{-2.6em}}
  {\titlerule*[0.7em]{\color{Rule}.}\sffamily\footnotesize\color{Muted}\contentspage}
\makeatletter
\renewcommand{\tableofcontents}{%
  \GuideOpenerTitle{}{\contentsname}%
  \GuideChapterMarkSet{Contents}\GuideLessonMarkSet{}%
  \setcounter{tocdepth}{0}% chapters and lessons only
  \@starttoc{toc}\clearpage}
\makeatother

% ── Openers ───────────────────────────────────────────────────────────────────
\newcounter{GuideChapter}
\newcounter{GuideLesson}[GuideChapter]
\newif\ifGuideInChapter
\newif\ifGuideMainMatter
\newcommand{\GuideChapterMarkSet}[1]{\renewcommand{\GuideChapterMark}{#1}}
\newcommand{\GuideLessonMarkSet}[1]{\renewcommand{\GuideLessonMark}{#1}}
\newcommand{\GuideStartMainMatter}{%
  \ifGuideMainMatter\else\GuideMainMattertrue\clearpage\pagenumbering{arabic}\pagestyle{guide}\fi}

% A lesson-style opener title: an optional small label, the title, an accent rule.
\newcommand{\GuideOpenerTitle}[2]{%
  \clearpage\thispagestyle{opener}%
  \vspace*{0.35in}%
  {\setlength{\parskip}{0pt}%
   \ifstrempty{#1}{}{{\HeadSemi\fontsize{12}{14}\selectfont\color{Accent}#1}\par\vspace{8pt}}%
   {\HeadFont\fontsize{23}{28}\selectfont\color{Ink}\raggedright #2\par}%
   \vspace{12pt}{\color{Accent}\rule{0.9in}{2.2pt}}\par}%
  \vspace{22pt}}

% Chapter opener: a Night band across the top of the page with the chapter
% number and title; the chapter introduction follows on the same page.
\newcommand{\GuideChapter}[3]{% number, title, anchor
  \GuideStartMainMatter
  \clearpage\thispagestyle{opener}%
  \GuideInChaptertrue
  \setcounter{GuideChapter}{#1}%
  % \setcounter doesn't reset dependent counters, so reset them here.
  \setcounter{GuideLesson}{0}\setcounter{GuideFigure}{0}%
  \GuideChapterMarkSet{Chapter #1 · #2}\GuideLessonMarkSet{Introduction}%
  \phantomsection\hypertarget{#3}{}\GuideSetLabel{#1}\label{#3}%
  \addcontentsline{toc}{part}{\protect\numberline{#1}#2}%
  \begin{tikzpicture}[remember picture, overlay]
    \fill[Night] (current page.north west) rectangle ([yshift=-3.05in]current page.north east);
    \begin{scope}[shift={([xshift=-5.2in, yshift=-3.05in]current page.north east)}]
      \clip (0,0) rectangle (5.2in, 3.05in);
      \GuideMotif{NightSoft}{0.4}{-0.4}{5}
    \end{scope}
    \fill[Accent] ([yshift=-3.05in]current page.north west) rectangle ([yshift=-3.12in]current page.north east);
    \node[anchor=south west, inner sep=0pt, align=left, text width=5.6in]
      at ([xshift=0.85in, yshift=-2.62in]current page.north west)
      {{\LabelFont\small\color{AccentBright}CHAPTER\par}\vspace{2pt}%
       {\DisplayFont\fontsize{58}{58}\selectfont\color{white}#1\par}\vspace{12pt}%
       {\HeadFont\fontsize{25}{30}\selectfont\color{white}\raggedright #2\par}};
  \end{tikzpicture}%
  \vspace*{2.35in}\par}

% Lesson opener: numbered N.k inside a chapter, unnumbered outside one
% (preface, reference pages).
\newcommand{\GuideLesson}[2]{% title, anchor
  \ifGuideInChapter
    \clearpage % ship the previous page before the number changes (running head)
    \stepcounter{GuideLesson}%
    \GuideOpenerTitle{\arabic{GuideChapter}.\arabic{GuideLesson}}{#1}%
    \GuideLessonMarkSet{\arabic{GuideChapter}.\arabic{GuideLesson}\enspace #1}%
  \else
    \GuideOpenerTitle{}{#1}%
    \GuideChapterMarkSet{#1}\GuideLessonMarkSet{}%
  \fi
  \vspace{-22pt}\phantomsection\hypertarget{#2}{}%
  \ifGuideInChapter\GuideSetLabel{\arabic{GuideChapter}.\arabic{GuideLesson}}\fi\label{#2}%
  \ifGuideInChapter
    \addcontentsline{toc}{chapter}{\protect\numberline{\arabic{GuideChapter}.\arabic{GuideLesson}}#1}%
  \else
    \addcontentsline{toc}{chapter}{#1}%
  \fi
  \vspace{22pt}}

% Chapter-end page (Exercises and Further Reading, Questions and Sources):
% inside the chapter, but unnumbered.
\newcommand{\GuideChapterEnd}[2]{% title, anchor
  \GuideOpenerTitle{}{#1}%
  \GuideLessonMarkSet{#1}%
  \vspace{-22pt}\phantomsection\hypertarget{#2}{}\label{#2}%
  \addcontentsline{toc}{chapter}{\protect\numberline{}#1}%
  \vspace{22pt}}
\makeatletter
\newcommand{\GuideSetLabel}[1]{\protected@edef\@currentlabel{#1}}
\makeatother

% Part page: a full Night page (used for "Reference").
\newcommand{\GuidePart}[2]{% title, anchor
  \GuideStartMainMatter
  \clearpage\thispagestyle{empty}\GuideInChapterfalse
  \phantomsection\hypertarget{#2}{}\label{#2}%
  \addcontentsline{toc}{part}{#1}%
  \begin{tikzpicture}[remember picture, overlay]
    \fill[Night] (current page.south west) rectangle (current page.north east);
    \begin{scope}[shift={(current page.south east)}, xshift=-11cm, yshift=-0.5cm]
      \clip (0,0) rectangle (11.5cm, 9cm);\GuideMotif{NightSoft}{0}{0}{5}
    \end{scope}
    \fill[AccentBright] ([xshift=0.85in, yshift=-3.3in]current page.north west) rectangle ++(0.9in, 4pt);
    \node[anchor=north west, inner sep=0pt]
      at ([xshift=0.85in, yshift=-3.55in]current page.north west)
      {\DisplayFont\fontsize{40}{44}\selectfont\color{white}#1};
  \end{tikzpicture}%
  \null\clearpage}

% Book parts: "Part II — Title" groups chapters. In a book that has them, the
% back-matter parts (Hints, Solutions, Reference) use the same page, unnumbered,
% so every part sits at the same level in the contents and the PDF outline.
\newcommand{\GuideTocPart}[3]{% numeral (may be empty), title, anchor
  \par\addvspace{1.5em}\Needspace*{5\baselineskip}%
  {\setlength{\parskip}{0pt}\noindent
   \ifstrempty{#1}{}{{\LabelFont\scriptsize\color{Accent}PART #1}\par\vspace{2pt}\noindent}%
   {\HeadFont\fontsize{12}{15}\selectfont\hyperlink{#3}{\color{Ink}#2}}\par
   \vspace{5pt}{\color{Rule}\hrule height 0.5pt}\par}%
  \vspace{0.35em}}
\newcommand{\GuideBookPart}[4]{% numeral (may be empty), title, anchor, bookmark text
  \GuideStartMainMatter
  \clearpage\thispagestyle{empty}\GuideInChapterfalse
  \phantomsection\hypertarget{#3}{}\label{#3}%
  \pdfbookmark[-2]{#4}{#3-part}%
  \addtocontents{toc}{\protect\GuideTocPart{#1}{#2}{#3}}%
  \GuideChapterMarkSet{#2}\GuideLessonMarkSet{}%
  \begin{tikzpicture}[remember picture, overlay]
    \fill[Night] (current page.south west) rectangle (current page.north east);
    \begin{scope}[shift={(current page.south east)}, xshift=-11cm, yshift=-0.5cm]
      \clip (0,0) rectangle (11.5cm, 9cm);\GuideMotif{NightSoft}{0}{0}{5}
      \GuideMotifAccent{AccentBright}{5/1, 6/1}
    \end{scope}
    \node[anchor=south west, inner sep=0pt, align=left, text width=5.6in]
      at ([xshift=0.85in, yshift=-3.25in]current page.north west)
      {\ifstrempty{#1}{}{{\LabelFont\small\color{AccentBright}PART\par}\vspace{2pt}%
        {\DisplayFont\fontsize{58}{58}\selectfont\color{white}#1\par}\vspace{14pt}}%
       {\color{AccentBright}\rule{0.9in}{4pt}}\par};
    \node[anchor=north west, inner sep=0pt, align=left, text width=5.6in]
      at ([xshift=0.85in, yshift=-3.55in]current page.north west)
      {\raggedright\DisplayFont\fontsize{34}{40}\selectfont\color{white}#2};
  \end{tikzpicture}%
  \null\clearpage}

% ── Problem sets ──────────────────────────────────────────────────────────────
% Problems at the end of each chapter; hints and solutions in their own parts at
% the back, so a reader never sees an answer by accident. Each problem, hint and
% solution links to the others by page number.
\newcommand{\GuideDifficulty}[2]{% level 1-3, label
  \tikz[baseline=-0.6ex]{\foreach \i in {1,2,3}{%
    \ifnum\i>#1\relax\fill[Rule] (\i*0.2,0) circle (0.06);%
    \else\fill[Accent] (\i*0.2,0) circle (0.06);\fi}}%
  \hspace{0.3em}{\LabelFont\scriptsize\color{Muted}\MakeUppercase{#2}}}
\newcommand{\GuideProblemRule}{{\color{Rule}\hrule height 0.6pt}\par\vspace{0.8em}}
\newcommand{\GuideProblem}[6]{% number, title, level, difficulty, target, signatures
  \par\addvspace{1.7em}\Needspace*{12\baselineskip}%
  \GuideProblemRule
  \phantomsection\hypertarget{problem-#1}{}\GuideSetLabel{#1}\label{problem-#1}%
  {\setlength{\parskip}{0pt}\noindent
   {\LabelFont\scriptsize\bfseries\color{AccentDark}PROBLEM #1}%
   \ifstrempty{#4}{}{\hspace{1em}\GuideDifficulty{#3}{#4}}\par
   \vspace{0.4em}{\HeadFont\fontsize{13.5}{17}\selectfont\color{Ink}\raggedright #2\par}%
   \ifstrempty{#6}{}{\vspace{0.35em}{\small\raggedright #6\par}}%
   \ifstrempty{#5}{}{\vspace{0.3em}{\LabelFont\tiny\color{Muted}TARGET}\hspace{0.5em}%
     {\sffamily\footnotesize\color{Ink}#5}\par}}%
  \vspace{0.25em}}
\newcommand{\GuideProblemLabel}[1]{%
  \par\Needspace*{4\baselineskip}\vspace{0.2em}{\LabelFont\scriptsize\bfseries\color{Muted}\MakeUppercase{#1}}\par\vspace{-0.3em}}
\newcommand{\GuideProblemEnd}[3]{% number, has hint (0/1), has solution (0/1)
  \par\vspace{0.1em}{\raggedleft\sffamily\scriptsize
   \ifnum#2=1 \hyperref[hint-#1]{Hint, page~\pageref*{hint-#1}}\fi
   \ifnum#2#3=11 {\color{Faint}\enspace·\enspace}\fi
   \ifnum#3=1 \hyperref[solution-#1]{Solution, page~\pageref*{solution-#1}}\fi\par}}
% The per-chapter pages inside the Hints and Solutions parts. Solutions start a
% new page per chapter; hints run on, since each chapter's hints are short.
\newcommand{\GuideAnswerPage}[5]{% part name, chapter number, chapter title, anchor, new page (0/1)
  \ifnum#5=1
    \GuideOpenerTitle{Chapter #2}{#3}\vspace{-22pt}%
  \else
    \par\addvspace{2em}\Needspace*{16\baselineskip}%
    {\setlength{\parskip}{0pt}%
     {\HeadSemi\fontsize{10}{12}\selectfont\color{Accent}Chapter #2}\par\vspace{4pt}%
     {\HeadFont\fontsize{17}{21}\selectfont\color{Ink}\raggedright #3\par}%
     \vspace{8pt}{\color{Accent}\rule{0.7in}{2pt}}\par}%
  \fi
  \phantomsection\hypertarget{#4}{}\label{#4}%
  \pdfbookmark[0]{Chapter #2 — #3}{#4-bm}%
  \GuideChapterMarkSet{#1}\GuideLessonMarkSet{Chapter #2\enspace #3}%
  \ifnum#5=1 \addcontentsline{toc}{chapter}{\protect\numberline{#2}#3}\fi
  \vspace{\ifnum#5=1 22pt\else 6pt\fi}}
\newcommand{\GuideHint}[2]{% number, title
  \par\addvspace{1.2em}\Needspace*{5\baselineskip}%
  \phantomsection\hypertarget{hint-#1}{}\GuideSetLabel{#1}\label{hint-#1}%
  {\setlength{\parskip}{0pt}\noindent
   {\LabelFont\scriptsize\bfseries\color{AccentDark}HINT #1}\hspace{0.8em}%
   {\HeadSemi\normalsize\color{Ink}#2}\hfill
   {\sffamily\scriptsize\hyperref[problem-#1]{Problem, page~\pageref*{problem-#1}}}\par}%
  \vspace{0.15em}}
\newcommand{\GuideSolution}[2]{% number, title
  \par\addvspace{1.9em}\Needspace*{10\baselineskip}%
  \GuideProblemRule
  \phantomsection\hypertarget{solution-#1}{}\GuideSetLabel{#1}\label{solution-#1}%
  {\setlength{\parskip}{0pt}\noindent
   {\LabelFont\scriptsize\bfseries\color{AccentDark}SOLUTION #1}\hfill
   {\sffamily\scriptsize\hyperref[problem-#1]{Problem, page~\pageref*{problem-#1}}}\par
   \vspace{0.4em}{\HeadFont\fontsize{13.5}{17}\selectfont\color{Ink}\raggedright #2\par}}%
  \vspace{0.3em}}
% A bold lead-in ("Brute force.") inside a solution.
\newcommand{\GuideRunIn}[1]{{\HeadSemi\color{AccentDark}#1}}

% ── Code panels ───────────────────────────────────────────────────────────────
% Long lines wrap with a marker instead of running off the page.
% pandoc defines Highlighting only when the book has highlighted code; a book
% with only plain blocks (guide.lua sets those in Highlighting too) needs it here.
\ifdefined\Highlighting\else\DefineVerbatimEnvironment{Highlighting}{Verbatim}{}\fi
\RecustomVerbatimEnvironment{Highlighting}{Verbatim}{%
  commandchars=\\\{\}, fontsize=\fontsize{8.4}{11.2}\selectfont,
  breaklines=true, breakanywhere=false,
  breaksymbolleft={\color{Faint}\tiny\ensuremath{\hookrightarrow}},
  breaksymbolindentleft=0.5em}
% A problem book keeps any listing that fits on a page in one piece
% (guide.lua sets \GuideKeepListingtrue around it).
\newif\ifGuideKeepListing
\tcbset{
  guide code/.style={enhanced, breakable, code={\ifGuideKeepListing\tcbset{unbreakable}\fi},
    colback=Panel, colframe=Panel, arc=2pt, boxrule=0pt,
    left=10pt, right=8pt, top=7pt, bottom=6pt, before skip=0.75em, after skip=0.9em,
    before={\par\needspace{4\baselineskip}}},
  guide code label/.style={overlay unbroken and first={%
    \ifstrempty{#1}{}{\node[anchor=north east, inner sep=3pt, font=\LabelFont\fontsize{6.5}{7}\selectfont, text=Muted]
      at ([xshift=-4pt, yshift=-2pt]frame.north east) {\MakeUppercase{#1}};}}},
  guide code title/.style={top=20pt, overlay unbroken and first app={%
    \node[anchor=north west, inner sep=0pt, font=\HeadSemi\fontsize{8}{9}\selectfont, text=Ink]
      at ([xshift=10pt, yshift=-7pt]frame.north west) {#1};
    \draw[Rule, line width=0.5pt] ([xshift=10pt, yshift=-15.5pt]frame.north west) -- ([xshift=-8pt, yshift=-15.5pt]frame.north east);}},
}
% Every code block sits on a warm panel with a thin accent edge and a small
% language label; a listing that is a whole file carries its file name as a title.
\newtcolorbox{GuideCode}[1]{guide code, borderline west={2pt}{0pt}{Accent!70},
  guide code label={#1}}
\newtcolorbox{GuideCodeTitled}[2]{guide code, borderline west={2pt}{0pt}{Accent!70},
  guide code label={#1}, guide code title={#2}}
% Terminal and plain-text blocks use the same panel with a neutral edge.
\newtcolorbox{GuideOutput}[1]{guide code, borderline west={2pt}{0pt}{Faint},
  guide code label={#1}}
\newtcolorbox{GuideOutputTitled}[2]{guide code, borderline west={2pt}{0pt}{Faint},
  guide code label={#1}, guide code title={#2}}

% ── Callouts ──────────────────────────────────────────────────────────────────
\newtcolorbox{GuideCallout}[3]{% ink colour, tint colour, label
  enhanced, breakable, colback=#2, colframe=#2, arc=2pt, boxrule=0pt,
  left=12pt, right=12pt, top=9pt, bottom=8pt, before skip=1em, after skip=1em,
  before={\par\needspace{6\baselineskip}},
  borderline west={3pt}{0pt}{#1},
  fontupper=\small\raggedright\setlength{\parskip}{0.45em},
  before upper={{\LabelFont\scriptsize\bfseries\color{#1}\MakeUppercase{#3}}\par\vspace{0.15em}}}
\newtcolorbox{GuideQuote}{% the Preface's opening passage
  enhanced, breakable, colback=white, colframe=white, boxrule=0pt, arc=0pt,
  left=14pt, right=0pt, top=2pt, bottom=2pt, before skip=0.4em, after skip=1.2em,
  borderline west={2.5pt}{0pt}{Accent},
  fontupper=\LiningSerif\fontsize{11.5}{17}\selectfont\color{Ink}\setlength{\parskip}{0.7em}}
% Example sentences in a language-learning book (a blockquote that starts with
% the target language): the sentence, its romanization, its translation.
\newtcolorbox{GuideExample}{%
  enhanced, breakable, colback=white, colframe=white, boxrule=0pt, arc=0pt,
  left=12pt, right=0pt, top=1pt, bottom=1pt, before skip=0.5em, after skip=0.7em,
  borderline west={1.5pt}{0pt}{Accent!60},
  fontupper=\setlength{\parskip}{0.3em}\color{Ink}}

% ── Sidebars ──────────────────────────────────────────────────────────────────
% "Going Deeper" digressions: optional reading, boxed so a reader can skip them.
\newtcolorbox{GuideSidebar}[1]{%
  enhanced, breakable, colback=Panel, colframe=Rule, arc=3pt, boxrule=0.6pt,
  left=14pt, right=14pt, top=11pt, bottom=10pt, before skip=1.3em, after skip=1.1em,
  before={\par\needspace{8\baselineskip}},
  fontupper=\small\setlength{\parskip}{0.45em},
  before upper={{\LabelFont\scriptsize\color{Muted}GOING DEEPER}\par\vspace{0.1em}%
    {\HeadSemi\normalsize\color{AccentDark}#1}\par\vspace{0.35em}}}

% ── Section summary ───────────────────────────────────────────────────────────
\newtcolorbox{GuideSummary}{%
  enhanced, breakable, code={\ifGuideKeepListing\tcbset{unbreakable}\fi}, colback=AccentTint, colframe=AccentTint, arc=3pt, boxrule=0pt,
  left=14pt, right=14pt, top=11pt, bottom=10pt, before skip=1.6em, after skip=0.6em,
  before={\par\needspace{8\baselineskip}},
  fontupper=\small\setlength{\parskip}{0.35em},
  before upper={{\HeadFont\normalsize\color{AccentDark}Summary}\par\vspace{0.3em}}}

% ── Figures ───────────────────────────────────────────────────────────────────
% A diagram and its caption never separate. Numbered per chapter (Figure 3.2).
\newcounter{GuideFigure}[GuideChapter]
\newenvironment{GuideFigure}
  {\par\vspace{0.6em}\noindent\begin{minipage}{\linewidth}\centering}
  {\end{minipage}\par\vspace{0.8em}}
\newcommand{\GuideCaption}[1]{%
  \par\vspace{0.55em}%
  \stepcounter{GuideFigure}%
  {\raggedright\small\setlength{\parskip}{0pt}%
   {\LabelFont\bfseries\scriptsize\color{AccentDark}FIGURE
    \ifGuideInChapter\arabic{GuideChapter}.\fi\arabic{GuideFigure}}\quad
   {\color{Muted}#1}\par}}

% ── Index ─────────────────────────────────────────────────────────────────────
% The heading is set by \GuideLesson, like every other page; imakeidx prints none.
\newcommand{\GuideNoHeading}[1]{}
\indexsetup{level=\GuideNoHeading, noclearpage, firstpagestyle=opener,
  othercode={\sffamily\footnotesize\setlength{\parskip}{0pt}}}
\newcommand{\GuideIndex}{\GuideInChapterfalse\GuideLesson{Index}{index}\printindex}

% ── Back cover ────────────────────────────────────────────────────────────────
\newcommand{\GuideBackCover}{%
  \clearpage\thispagestyle{empty}%
  \begin{tikzpicture}[remember picture, overlay]
    \fill[Night] (current page.south west) rectangle (current page.north east);
    \begin{scope}[shift={(current page.south east)}, xshift=-11cm, yshift=-0.5cm]
      \clip (0,0) rectangle (11.5cm, 5.5cm);\GuideMotif{NightSoft}{0}{0}{3}
    \end{scope}
    \node[anchor=north west, text width=5.5in, inner sep=0pt, align=left]
      at ([xshift=0.95in, yshift=-1.1in]current page.north west)
      {\setlength{\parskip}{0.9em}\raggedright
       {\LabelFont\small\color{AccentBright}\MakeUppercase{\GuideCategory}\par}
       \vspace{0.2em}{\color{AccentBright}\rule{0.9in}{3pt}}\par\vspace{0.4em}
       {\LiningSerif\fontsize{12}{18}\selectfont\color{white}\GuideBlurb\par}};
    \node[anchor=south west, inner sep=0pt, align=left, text width=4.5in]
      at ([xshift=0.95in, yshift=0.8in]current page.south west)
      {{\HeadFont\normalsize\spaceskip=0.26em\relax\color{white}\GuideTitle}\par\vspace{3pt}
       {\LabelFont\scriptsize\color{Faint}\MakeUppercase{\GuideCoverage}}};
  \end{tikzpicture}}

% ── Tables ────────────────────────────────────────────────────────────────────
\newcommand{\GuideTableSetup}{%
  \sffamily\small\setlength{\tabcolsep}{8pt}\renewcommand{\arraystretch}{1.45}%
  \arrayrulecolor{Rule}\rowcolors{2}{white}{Panel}}
\let\GuideOrigToprule\toprule
\renewcommand{\toprule}{\arrayrulecolor{Accent}\specialrule{1pt}{0pt}{0pt}\arrayrulecolor{Rule}}
\renewcommand{\bottomrule}{\arrayrulecolor{Rule}\specialrule{0.8pt}{0pt}{0pt}}

% ── Glossary ──────────────────────────────────────────────────────────────────
\newcommand{\GuideGlossaryLetter}[1]{%
  \par\vspace{1em}\needspace{4\baselineskip}%
  {\DisplayFont\fontsize{18}{20}\selectfont\color{Accent}#1}\par
  {\color{Rule}\rule{\linewidth}{0.5pt}}\par\vspace{0.2em}}
